\uuid{N8Hv}
\exo7id{4807}
\titre{exo7 4807}
\auteur{quercia}
\organisation{exo7}
\datecreate{2010-03-16}
\isIndication{false}
\isCorrection{false}
\chapitre{Topologie}
\sousChapitre{Topologie des espaces vectoriels normés}
\cpge{mpsi-mp2i.applications-lineaires.c, mpsi-mp2i.matrices-applications-lineaires.a, pcsi.matrices-applications-lineaires.a, ptsi.matrices-applications-lineaires.a, bcpst1.applications-lineaires-matrices, mp-mpi.topologie-espaces-normes.e, pc.espaces-vectoriels-normes.e, psi.espaces-vectoriels-normes.e, pt.fonctions-deux-trois-variables.b, mp-mpi.topologie-espaces-normes.f, mp-mpi.topologie-espaces-normes.j, pc.espaces-vectoriels-normes.g, psi.espaces-vectoriels-normes.g, bcpst2.applications-lineaires-matrices.a}
\module{Analyse}
\niveau{L2}
\difficulte{}

\contenu{
\texte{
Soient $x_1,\dots,x_n \in \R$ distincts.
{\`A} $(y_1,\dots,y_n) \in \R^n$ on fait correspondre le polyn{\^o}me
$P \in \R_{n-1}[X]$ tel que pour tout $i$ : $P(x_i) = y_i$.
}
\begin{enumerate}
    \item \question{Montrer que l'application $(y_1,\dots,y_n)  \mapsto P$ est continue.}
    \item \question{Montrer que l'application r{\'e}ciproque est aussi continue.}
\end{enumerate}
}
